\documentclass[12pt]{article}
\usepackage[T1]{fontenc}
\usepackage[utf8]{inputenc}
\usepackage{lmodern}
\usepackage{textcomp}
\usepackage{enumitem}
\usepackage{geometry}
\geometry{margin=1in}
\begin{document}
\begin{center}
Answer Key - Physics - Graduate Level Assessment 24 \
\small (Answers are ordered exactly as in the question paper.)
\end{center}

\section*{Multiple Choice}
\begin{enumerate}[label=\arabic*.]
\item \textbf{Answer:} A
\\ \textit{Options:} A) 33.9; B) 29.8; C) 27.5; D) 31.2; E) 24.3
\\ \textit{Note:} The binding energy of the 4 He nucleus, calculated using the mass defect and Einstein's equation E=mc², yields approximately 33.9 MeV, reflecting the energy required to hold the nucleus together against the repulsive forces between protons.
\item \textbf{Answer:} A
\\ \textit{Options:} A) 0.4 henry; B) 0.3 henry; C) 0.2 henry; D) 0.6 henry; E) 0.7 henry
\\ \textit{Note:} The self-inductance of the coil can be calculated using the formula \( L = \frac{V}{\frac{di}{dt}} \), where \( V \) is the back electromotive force (2 volts) and \( \frac{di}{dt} \) is the rate of change of current (5 amperes per second), resulting in \( L = \frac{2 \, \text{V}}{5 \, \text{A/s}} = 0.4 \, \text{H} \).
\item \textbf{Answer:} B
\\ \textit{Options:} A) It is increased by two times.; B) It remains the same.; C) It becomes zero.; D) It is reduced to half.; E) It is increased by eight times.
\\ \textit{Note:} The gravitational force between two objects is directly proportional to the product of their masses and inversely proportional to the square of the distance between them, so doubling both the masses and the distance results in the force remaining unchanged, as both changes cancel each other out.
\item \textbf{Answer:} E
\\ \textit{Options:} A) (f / 3); B) (f / 2); C) (5 f / 3); D) (6 f / 5); E) (3 f / 4)
\\ \textit{Note:} The correct answer is derived from the lens formula for thin lenses in combination, where the effective focal length (F) of two lenses separated by a distance (d) is given by the formula F = (f 1 * f 2) / (f 1 + f 2 - d) for lenses with focal lengths f 1 and f 2, leading to F = (f * f) / (f + f - (2 f/3)), which simplifies to (3 f/4).
\item \textbf{Answer:} A
\\ \textit{Options:} A) because Mars' axis is less tilted than the Earth's; B) because Mars has a smaller size than the Earth; C) because Mars has a slower rotation speed than the Earth; D) because Mars has more carbon dioxide in its atmosphere than the Earth; E) because Mars has a shorter year than the Earth
\\ \textit{Note:} The seasons in Mars' southern hemisphere are more extreme than on Earth because Mars has a smaller axial tilt (approximately 25 degrees compared to Earth's 23.5 degrees), resulting in greater variations in solar radiation received during its orbital path, which amplifies seasonal changes.
\end{enumerate}

\section*{True / False}
\begin{enumerate}[label=\arabic*.]
\setcounter{enumi}{5}
\item \textbf{Answer:} True
\\ \textit{Options:} A) True; B) False
\\ \textit{Note:} The statement is true because a ratio of forward-bias to reverse-bias currents of -6 e+25 indicates that the forward current is vastly greater than the reverse current, reflecting the exponential increase in carrier injection under forward-bias conditions compared to the negligible reverse current in a diode.
\item \textbf{Answer:} False
\\ \textit{Options:} A) True; B) False
\\ \textit{Note:} The statement is false because an object with an angular size of 0.5 degrees and a height of 5 cm would actually be approximately 6 meters away from the observer, not 15 meters, based on the relationship between angular size, actual size, and distance.
\item \textbf{Answer:} False
\\ \textit{Options:} A) True; B) False
\\ \textit{Note:} A lens with a focal length of 20 inches cannot produce a virtual image when the object is placed 10 inches from it because the object is located within the focal length, which results in a real image rather than a virtual image.
\item \textbf{Answer:} False
\\ \textit{Options:} A) True; B) False
\\ \textit{Note:} The statement is false because a thin film's absorption is not solely determined by its optical density but also by its thickness and refractive index, which together influence the fraction of incident light that is absorbed.
\item \textbf{Answer:} True
\\ \textit{Options:} A) True; B) False
\\ \textit{Note:} The statement is true because the Sun's surface temperature of approximately 6000 K places it in the G-type star category, which aligns with the established stellar temperature classification system that categorizes stars based on their effective temperatures.
\end{enumerate}

\section*{Long Form Response}
\begin{enumerate}[label=\arabic*.]
\setcounter{enumi}{10}
\item \textbf{Answer:} When 100 g of water at 0^\ensuremath{\circC} is mixed with 50 g of water at 50^\ensuremath{\circC}, the system reaches thermal equilibrium, resulting in a change in entropy. The final temperature can be calculated using the principle of conservation of energy, leading to a mixed temperature of approximately 20^\ensuremath{\circC}. The change in entropy (ΔS) for the system can be determined using the formula ΔS = mc ln(T\_final/T\_initial), where m is the mass, c is the specific heat capacity of water (1 cal/g·^\ensuremath{\circC}), and T is the temperature in Kelvin. Considering the contributions from both masses of water and the calculated final temperature, the total change in entropy for the system is approximately 0.5 cal/K/deg. This positive change reflects the increase in disorder as the two water masses mix, consistent with the second law of thermodynamics.
\\ \textit{Note:} correct because it correctly identifies the process of reaching thermal equilibrium, applies the conservation of energy to find the final temperature, and uses the appropriate entropy formula to quantify the change in entropy for the mixed water system.
\item \textbf{Answer:} To calculate the electrostatic force acting on a 10.0 g block with a charge of +8.00 \ensuremath{\times} 10^-5 C in the electric field E = (3000 i - 600 j) N/C, we use the equation F = qE, where F is the force, q is the charge, and E is the electric field vector. Substituting the given values, we find F = (8.00 \ensuremath{\times} 10^-5 C) * (3000 i - 600 j) N/C, which results in a force vector of (0.24 i - 0.048 j) N. The magnitude of this force can be calculated using the Pythagorean theorem, yielding approximately 0.245 N. This calculation illustrates the fundamental principle that the electrostatic force on a charge in an electric field is directly proportional to the magnitude of the charge and the strength of the electric field, highlighting the interaction between electric charges and fields in electrostatics.
\\ \textit{Note:} correct because it accurately applies the formula for electrostatic force, F = qE, using the given charge and electric field vector, and correctly computes both the vector and magnitude of the resulting force.
\end{enumerate}

\vfill
\noindent\textit{End of Answer Key}
\end{document}